% Upload this file (or the whole docs/overleaf/ folder) to Overleaf.
% Compile with pdfLaTeX.
% Companion to rff_woodbury_derivation.tex (scalar RFF Woodbury).
\documentclass[11pt]{article}

\usepackage[margin=1in]{geometry}
\usepackage{amsmath,amssymb,mathtools}
\usepackage{booktabs}
\usepackage{colortbl}
\usepackage{xcolor}
\usepackage{enumitem}
\usepackage{tcolorbox}
\usepackage{hyperref}
\usepackage{natbib}

\hypersetup{colorlinks=true, linkcolor=blue!60!black, citecolor=blue!60!black}

\tcbset{enhanced, boxrule=0.6pt, arc=2pt, left=8pt, right=8pt, top=6pt, bottom=6pt}

\newtcolorbox{samebox}{
  colback=green!6,
  colframe=green!50!black,
  title={\textbf{Reuses scalar Woodbury} (see \texttt{rff\_woodbury\_derivation.tex})},
}

\newtcolorbox{diffbox}{
  colback=orange!8,
  colframe=orange!70!black,
  title={\textbf{New for multitask}},
}

\newtcolorbox{planbox}{
  colback=blue!5,
  colframe=blue!50!black,
  title={\textbf{Implementation action}},
}

\newcommand{\noise}{\sigma^2}
\newcommand{\Sig}{\Sigma}
\newcommand{\Ph}{\Phi}
\newcommand{\Om}{\Omega}
\newcommand{\LB}{L_B}
\newcommand{\RB}{R_B}

\title{Multitask RFF-GP Plan:\\[0.3em]
Combining \texttt{MTGPR} with ORF/RFF/SORF Woodbury Inference}
\author{GPPlus --- implementation plan}
\date{\today}

\begin{document}
\maketitle

\begin{abstract}
This document is an \textbf{implementation plan} for joint multitask Gaussian process regression
with Random Fourier Features (RFF), including Orthogonal RFF (ORF) and Structured ORF (SORF),
at Woodbury cost instead of dense $O(n^3 T^3)$ exact inference.
Today GPPlus has two disconnected stacks: \texttt{MTGPR} (multitask targets, dense kernel) and
\texttt{RFFGPR} (scalar targets, Woodbury). The plan generalizes scalar Woodbury
($\Sig = \noise I_n + \Ph\Ph^\top$) to the intrinsic coregionalization model (ICM)
Kronecker form $\Sig = K_x \otimes B + \Lambda$ via a joint feature matrix
$\Om = \Ph \otimes \RB$ via a PSD factor of the full task matrix $B$.
Green boxes mark mathematics that directly extends the scalar note;
orange boxes mark new multitask-specific design; blue boxes mark concrete code tasks.
Section~\ref{sec:mtgpr} documents \textbf{what \texttt{MTGPR} does} in GPPlus/GPyTorch---the dense
reference whose multitask covariance the proposed Woodbury path must reproduce.
Section~\ref{sec:related} situates the plan in the multitask-GP and RFF literature.
\end{abstract}

\tableofcontents
\newpage

%=============================================================================
\section{Goal and motivation}
%=============================================================================

\subsection{What we want}

Train a \textbf{single} Gaussian process on multitask data
$y \in \mathbb{R}^{n \times T}$ (e.g.\ TOA: $y_{\cos}$, $y_{\mathrm{grain}}$) with:
\begin{itemize}[nosep]
  \item shared spatial inputs $X \in \mathbb{R}^{n \times d}$;
  \item \textbf{learned cross-task correlation} (ICM / \texttt{MultitaskKernel});
  \item \textbf{RFF/ORF/SORF} feature approximation for scalability in $n$ and $d$;
  \item \textbf{Woodbury} training and prediction ($O((mT)^3)$ when using a full $T\times T$ task factor; not $O(n^3 T^3)$).
\end{itemize}

\subsection{What exists today}

\begin{center}
\renewcommand{\arraystretch}{1.3}
\begin{tabular}{@{}p{0.22\textwidth}p{0.34\textwidth}p{0.34\textwidth}@{}}
\toprule
\textbf{Component} & \textbf{\texttt{MTGPR}} & \textbf{\texttt{RFFGPR}} \\
\midrule
Targets & $(n, T)$ multitask & scalar $(n,)$ \\
Kernel & \texttt{MultitaskKernel(GaussianKernel)} & \texttt{LogScaleKernel(RFFKernel)} \\
Likelihood & \texttt{MultitaskGaussianLikelihood} & \texttt{LogGaussianLikelihood} \\
Inference & dense \texttt{ExactMarginalLogLikelihood} & \texttt{RFFWoodburyMarginalLogLikelihood} \\
Cost (rough) & $O(n^3 T^3)$ & $O(n m^2)$, $m = 2D$ \\
ORF/RFF/SORF & no & yes (\texttt{rff\_sampling}) \\
Experiments & \texttt{MT1\_analytic\_1d\_GP.py} & \texttt{S1\_toa\_*}, 2$\times$ scalar loop \\
\bottomrule
\end{tabular}
\end{center}

\begin{diffbox}
\textbf{Cannot merge by wrapping \texttt{RFFKernel} inside existing \texttt{MTGPR} alone.}
GPyTorch's exact MLL still materializes a dense multitask covariance.
The current TOA workaround (two independent \texttt{RFFGPR} models) does \emph{not} learn task correlation.
\end{diffbox}

%=============================================================================
\section{Related work}
\label{sec:related}
%=============================================================================

This plan combines three well-studied ideas---multitask ICM GPs, random Fourier features, and
Woodbury low-rank inference---into a single GPPlus implementation path.
No one reference does exactly \texttt{RFFMTGPR} (GPyTorch \texttt{MultitaskKernel} + ORF/RFF/SORF +
$\Om = \Ph \otimes \RB$ Woodbury), but each ingredient and the natural combination are grounded
in prior work.

\subsection{Multitask and coregionalized GPs}

The probabilistic target is the intrinsic coregionalization model (ICM): a learnable task
covariance $B$ multiplies a shared spatial kernel,
$\mathrm{Cov}[f_t(x), f_{t'}(x')] = B_{t,t'}\, k(x,x')$,
giving joint covariance $K_x \otimes B$ under the usual vec ordering
(Section~\ref{sec:mtgpr}).
\citet{bonilla2008mtgp} introduced this \emph{free-form} task matrix for multi-task GP prediction
and noted that exact inference requires $O(n^3 T^3)$ linear algebra on the $nT \times nT$
Gram matrix; their Section~3 uses Nystr\"om approximation of $K_x$ and a low-rank PCA of $B$
rather than RFF features with Woodbury.
\citet{alvarez2012kernels} review the linear model of coregionalization (LCM) and place ICM as
a central special case for vector-valued GPs.
GPyTorch's \texttt{MultitaskKernel}~\citep{gardner2018gpytorch} implements the same
$K_{XX} \otimes K_{TT}$ structure and is the dense reference for \texttt{MTGPR}.

\subsection{Scalable Kronecker inference (without RFF)}

When inputs lie on grids or admit Kronecker structure, exact GPs can avoid dense $nT$ Cholesky:
\citet{wilson2015kronecker} extend Kronecker methods to non-Gaussian likelihoods via conjugate
gradients; \citet{lin2025latentkronecker} use latent Kronecker structure with iterative solvers
for large $n$.
\citet{leroy2023lmc} give decoupled exact inference for the full LMC without restrictive
orthogonality assumptions on the mixing matrix.
These methods exploit \emph{exact} $K_x$, whereas \texttt{RFFMTGPR} deliberately approximates
$K_x \approx \Ph\Ph^\top$ and targets the same ICM prior at Woodbury cost when $mT \ll nT$.

\subsection{Random Fourier features and Woodbury for GPs}

\citet{rahimi2007rff} introduced RFF to scale kernel machines by replacing implicit kernel
matrices with finite feature maps $\Ph$.
For GP regression, a low-rank covariance $\Ph\Ph^\top + \noise I_n$ invites the Woodbury
identity, reducing inference from $O(n^3)$ to $O(m^3)$ when $m \ll n$; GPPlus scalar
\texttt{RFFGPR} follows this template (see \texttt{rff\_woodbury\_derivation.tex}).
\citet{cutajar2017dgp} formulate scalable deep GPs via random feature expansions and Woodbury
marginal likelihood (with stochastic variational inference); Bonilla is a co-author, linking
this line to multi-task GP research, though the published focus is deep/single-task structure
rather than ICM Kronecker features.
\citet{potapczynski2021rfftrunc} discuss bias from truncating RFF in GP hyperparameter learning---a
caveat that applies equally to multitask RFF.

\subsection{Multitask learning with shared RFFs}

\citet{sevilla2023rffblr} propose RFF-BLR: shared random Fourier features across tasks within a
Bayesian linear regression framework, promoting sparsity and transfer without a full GP
marginal likelihood or ICM Woodbury factorization.
That work is closest in \emph{motivation} (multitask + RFF) but differs in model class and
inference; \texttt{RFFMTGPR} aims for parity with dense \texttt{MTGPR} / \texttt{MultitaskKernel},
not a sparse ELM-style regressor.

\subsection{Positioning of the GPPlus plan}

\begin{center}
\renewcommand{\arraystretch}{1.25}
\begin{tabular}{@{}p{0.30\textwidth}p{0.62\textwidth}@{}}
\toprule
\textbf{Prior work} & \textbf{Relation to \texttt{RFFMTGPR}} \\
\midrule
Bonilla et al.\ ICM MTGP & Same generative model; they approximate $K_x$ and $B$ for speed \\
Cutajar et al.\ RFF + Woodbury & Same inference trick as scalar \texttt{RFFGPR}; not multitask \\
Kronecker / LMC fast inference & Exact $K_x$ structure; no RFF approximation \\
RFF-BLR & Multitask + shared RFF; not ICM GP with exact Woodbury MLL \\
\textbf{GPPlus plan} & ICM prior + $\Om=\Ph\otimes\RB$ + Woodbury + ORF/SORF \\
\bottomrule
\end{tabular}
\end{center}

The engineering contribution is to \textbf{unify} GPPlus's existing stacks---dense \texttt{MTGPR}
(reference ICM) and Woodbury \texttt{RFFGPR} (reference scalable inference)---with
implementation details (full $B$, vec order, per-task $\Lambda$) verified against GPyTorch.

%=============================================================================
\section{What \texttt{MTGPR} does (dense reference model)}
\label{sec:mtgpr}
%=============================================================================

This section records the \textbf{exact multitask mathematics} implemented by
\texttt{gpplus/models/mtgpr.py} and GPyTorch's exact-GP stack.
The proposed \texttt{RFFMTGPR} must match this object (same prior/likelihood structure;
only replace dense $k(x,x')$ by $\Ph\Ph^\top$ and use Woodbury instead of $O(n^3T^3)$ Cholesky).

\subsection{Generative model}

At each of $n$ input locations $x_i \in \mathbb{R}^d$, define a latent task vector
$f(x_i) \in \mathbb{R}^T$. Stacked observations are
\begin{equation}
  Y \in \mathbb{R}^{n \times T}, \qquad Y_{i,t} = f_t(x_i) + \varepsilon_{i,t}.
\end{equation}
GPPlus stores $Y$ as \texttt{train\_y} with shape \texttt{(n, num\_tasks)}.
Each \textbf{row} is one spatial location; each \textbf{column} is one task.

\subsection{Three modules in \texttt{MTGPR}}

\texttt{MTGPR} is a thin wrapper around \texttt{gpytorch.models.ExactGP} with three defaults
(see \texttt{mtgpr.py}):

\begin{center}
\renewcommand{\arraystretch}{1.25}
\begin{tabular}{@{}p{0.28\textwidth}p{0.62\textwidth}@{}}
\toprule
\textbf{Module} & \textbf{Role} \\
\midrule
\texttt{MultitaskMean} & Wraps a scalar mean (default \texttt{ConstantMean}) to produce
$\mu \in \mathbb{R}^{n \times T}$ with one constant per task. \\
\texttt{MultitaskKernel} & ICM kernel: shared spatial $k(x,x')$ times learnable task matrix $B$. \\
\texttt{MultitaskGaussianLikelihood} & Gaussian noise on $Y$; per-task diagonal if
\texttt{rank\_likelihood}$=0$, or low-rank task noise if $>0$. \\
\bottomrule
\end{tabular}
\end{center}

Hyperparameters \texttt{rank\_kernel} ($r_k$) and \texttt{rank\_likelihood} ($r_\ell$)
control the rank of $B$ and of the noise covariance, respectively.

\subsection{Forward pass}

For inputs $X \in \mathbb{R}^{n \times d}$:
\begin{equation}
  \texttt{model(X)} \;\Rightarrow\;
  \mathcal{N}\bigl(\mu(X),\; \mathcal{K}_{\mathrm{MT}}(X)\bigr),
\end{equation}
a \texttt{MultitaskMultivariateNormal} with mean $\mu \in \mathbb{R}^{n \times T}$ and
a lazy $(nT) \times (nT)$ covariance $\mathcal{K}_{\mathrm{MT}}$ built from
\texttt{covar\_module} (the \texttt{MultitaskKernel}).

Training uses \texttt{ExactMarginalLogLikelihood}, which materializes this covariance
(or its Cholesky factor) and evaluates
\begin{equation}
  \log p(Y \mid X)
  = -\tfrac{1}{2}\,\mathrm{vec}(\tilde{Y})^\top \Sig^{-1}\,\mathrm{vec}(\tilde{Y})
  - \tfrac{1}{2}\log|\Sig| + \mathrm{const},
\end{equation}
with $\tilde{Y} = Y - \mu(X)$ and $\Sig$ the \textbf{joint} $(nT)\times(nT)$ training covariance
(spatial kernel $\otimes$ task kernel, plus noise).

\subsection{ICM kernel (what \texttt{MultitaskKernel} encodes)}

For spatial base kernel $k(\cdot,\cdot)$ (default \texttt{LogScaleKernel(GaussianKernel)},
or any wrapped base such as \texttt{GaussianKernel * PeriodicKernel} in \texttt{MT1\_analytic\_1d\_GP.py}):
\begin{equation}
  \label{eq:icm-element}
  \mathrm{Cov}\bigl[f_t(x_i), f_{t'}(x_j)\bigr] = B_{t,t'}\, k(x_i, x_j),
  \qquad
  B \in \mathbb{R}^{T \times T}\ \text{PSD (full task covariance matrix)}.
\end{equation}
GPyTorch's \texttt{IndexKernel} (\texttt{task\_covar\_module}) parameterizes $B$ as a
low-rank term plus a learned diagonal,
$B = \LB\LB^\top + \mathrm{diag}(\mathbf{v})$ with $\LB \in \mathbb{R}^{T \times r_k}$,
\textbf{not} simply $B = \LB\LB^\top$.
For Woodbury, use the dense matrix
\texttt{task\_covar\_module.covar\_matrix.to\_dense()} (or an equivalent PSD factor
$\RB \in \mathbb{R}^{T \times T}$ with $B = \RB\RB^\top$).
When \texttt{rank\_kernel}$=0$, $B$ is diagonal (independent tasks in the kernel).
Let $K_x \in \mathbb{R}^{n \times n}$ with $(K_x)_{ij} = k(x_i, x_j)$.
Then the \textbf{latent} covariance over all $nT$ entries of $Y$ is the Kronecker product
\begin{equation}
  \label{eq:kron-gpytorch}
  \boxed{
    \Sig_{\mathrm{latent}}
    = K_x \otimes B
    \;\in\; \mathbb{R}^{(nT)\times(nT)}.
  }
\end{equation}

\begin{samebox}
\textbf{GPyTorch convention (verified).}
Index observations as $p = t + T(i-1)$ for row $i$ and task $t$ (tasks vary \textbf{fast} within
each spatial point). With this ordering,
$\Sig_{\mathrm{latent}} = K_x \otimes B$ matches \texttt{MultitaskKernel.evaluate()}.
Equivalently, block $(i,j)$ of $\Sig_{\mathrm{latent}}$ is the $T\times T$ matrix $k(x_i,x_j)\,B$.
\end{samebox}

\begin{diffbox}
\textbf{Common pitfall (Kronecker order).} $K_x \otimes B$ is \emph{not} the same as $B \otimes K_x$.
The joint feature identity below uses $\Ph \otimes \RB$, not $\RB \otimes \Ph$.
Always regression-test vector ordering against dense \texttt{MTGPR}.

\textbf{Common pitfall (task factor).} \texttt{covar\_factor} alone does \emph{not} satisfy
$\LB\LB^\top = B$ in GPyTorch; do not set $\Om = \Ph \otimes \LB$ unless you have verified
$\LB\LB^\top = B$ numerically. Use a PSD factor of the full $B$ (Cholesky or symmetric
eigendecomposition square root).
\end{diffbox}

\subsection{Likelihood (what \texttt{MultitaskGaussianLikelihood} adds)}

Observations add noise on top of the latent:
\begin{equation}
  \Sig = \Sig_{\mathrm{latent}} + \Lambda,
  \qquad
  \Lambda = \text{noise from \texttt{MultitaskGaussianLikelihood}}.
\end{equation}
\textbf{Phase~1 (MVP):} \texttt{rank\_likelihood}$=0$ gives independent task noise variances
$\Lambda = \mathrm{diag}(\noise_1,\ldots,\noise_T)$ repeated across spatial blocks
(equivalently $\Lambda = I_n \otimes D_\noise$ with $D_\noise = \mathrm{diag}(\noise_t)$).
\textbf{Phase~2:} $r_\ell > 0$ adds correlated task noise; Woodbury must whiten or extend $\Lambda$.

\subsection{Prediction}

At test inputs $X_\ast \in \mathbb{R}^{n_\ast \times d}$, \texttt{MTGPR} returns
$\mathbb{E}[Y_\ast \mid X, Y] \in \mathbb{R}^{n_\ast \times T}$ and task correlations
within each spatial block (and across space via the GP). This is the target behavior for
\texttt{RFFMTGPR.predict}: same $(n_\ast, T)$ mean shape; per-task marginal std optional in MVP,
full $T\times T$ predictive blocks in Phase~2.

\subsection{Why \texttt{MTGPR} is the math reference}

\begin{enumerate}[nosep]
  \item It is the \textbf{only} GPPlus model today that trains on $(n,T)$ targets jointly.
  \item Its $\Sig = K_x \otimes B + \Lambda$ is the object we approximate with
        $K_x \approx \Ph\Ph^\top$ and Woodbury on $\Om\Om^\top + \Lambda$.
  \item Swapping \texttt{GaussianKernel} for \texttt{RFFKernel} inside \texttt{MultitaskKernel}
        without a custom MLL still triggers dense $(nT)^3$ algebra---hence a new model class.
  \item Any \texttt{RFFMTGPR} implementation should pass:
        \emph{``NLL and predictions match \texttt{MTGPR} at small $n,T$ when
        $k \approx \Ph\Ph^\top$''} (same full $B$, same noise, same vec order).
\end{enumerate}

%=============================================================================
\section{ICM multitask kernel and RFF substitution}
%=============================================================================

\subsection{Notation}

\begin{itemize}[nosep]
  \item $n$ training points, $T$ tasks, input dimension $d$.
  \item $D$ RFF frequencies; feature dimension $m = 2D$ (cos/sin pairs).
  \item $r_k$: rank of task kernel (\texttt{rank\_kernel} in \texttt{MultitaskKernel}).
  \item $r_\ell$: rank of task noise (\texttt{rank\_likelihood} in \texttt{MultitaskGaussianLikelihood}).
\end{itemize}

\subsection{Spatial RFF features}

For each input $x_i$, GPPlus builds scaled RFF features (same as scalar \texttt{RFFGPR}):
\begin{equation}
  \label{eq:phi}
  \Ph \in \mathbb{R}^{n \times m},
  \qquad
  \Ph_{i,:} = \mathbf{z}(x_i)^\top,
\end{equation}
with $\mathbf{z}(x) \in \mathbb{R}^{m}$ from \texttt{featurize\_rbf} and output-scale applied in
\texttt{RFFGPR.scaled\_features}.
ORF and SORF only change \textbf{how} frequency weights are drawn (\texttt{rff\_sampling});
Woodbury structure is unchanged (see \texttt{rff\_woodbury\_derivation.tex}, Section~ORF).

\subsection{Task covariance (ICM, rank $r_k$)}

GPyTorch's \texttt{MultitaskKernel} uses intrinsic coregionalization:
\begin{equation}
  \label{eq:icm}
  K\bigl((x_i, t), (x_j, t')\bigr) = B_{t,t'}\, k(x_i, x_j),
  \qquad
  B \in \mathbb{R}^{T \times T}\ \text{PSD}
  \quad\text{(from \texttt{covar\_matrix}; see Section~\ref{sec:mtgpr}).}
\end{equation}
With RFF approximation $K_x \approx \Ph\Ph^\top$,
the \textbf{full} latent training covariance (before noise) is
\begin{equation}
  \label{eq:sigma-mt-latent}
  \Sig_{\mathrm{latent}} = K_x \otimes B
  \;\approx\; (\Ph\Ph^\top) \otimes B,
\end{equation}
and the \textbf{observation} covariance is
\begin{equation}
  \label{eq:sigma-mt}
  \boxed{
    \Sig = K_x \otimes B + \Lambda
    \;\approx\; (\Ph\Ph^\top) \otimes B + \Lambda,
  }
\end{equation}
where $\Lambda$ is from \texttt{MultitaskGaussianLikelihood}
(diagonal per-task when $r_\ell = 0$). See Section~\ref{sec:mtgpr} for vec ordering.

\subsection{Vectorization order}

Let $Y \in \mathbb{R}^{n \times T}$ with entries $Y_{i,t}$ (\texttt{train\_y} layout).
GPyTorch uses \textbf{row-wise} stacking into $\mathrm{vec}(Y) \in \mathbb{R}^{nT}$:
\begin{equation}
  \mathrm{vec}(Y)_{\,t + T(i-1)} = Y_{i,t},
  \qquad
  t = 1,\ldots,T,\;\; i = 1,\ldots,n,
\end{equation}
so task index $t$ varies \textbf{fastest} and spatial index $i$ varies slowest.
Under this convention $\Sig_{\mathrm{latent}} = K_x \otimes B$ (Equation~\eqref{eq:kron-gpytorch}).
\textbf{All Kronecker constructions must use this ordering} (regression test vs dense \texttt{MTGPR}).

%=============================================================================
\section{Kronecker Woodbury factorization}
%=============================================================================

\subsection{Key identity}

Let $\RB \in \mathbb{R}^{T \times T}$ satisfy $B = \RB\RB^\top$ (any PSD square root of the
\textbf{full} task matrix $B$).
For $K_x = \Ph\Ph^\top$ and GPyTorch vec order (Section~\ref{sec:mtgpr}):
\begin{equation}
  \label{eq:kronecker-gram}
  K_x \otimes B
  = (\Ph \otimes \RB)\,(\Ph \otimes \RB)^\top.
\end{equation}

Define the \textbf{joint feature matrix} (note order: spatial $\otimes$ task)
\begin{equation}
  \label{eq:omega}
  \boxed{
    \Om = \Ph \otimes \RB
    \;\in\; \mathbb{R}^{(nT) \times (mT)}.
  }
\end{equation}
Then $\Om\Om^\top = K_x \otimes B = \Sig_{\mathrm{latent}}$ (before adding $\Lambda$).
Columns are $mT$ in general (not $m r_k$), because $B$ is $T \times T$ even when
\texttt{rank\_kernel} is small.

\begin{diffbox}
\textbf{Not} $\Om = \RB \otimes \Ph$. That factorization yields $B \otimes K_x$, which is the
wrong Kronecker order for GPyTorch \texttt{MTGPR}. Implementation and unit tests must use
\texttt{torch.kron(Phi, R\_B)} with $\Ph \in \mathbb{R}^{n \times m}$, $\RB \in \mathbb{R}^{T \times T}$
(\texttt{.contiguous()} on both factors before \texttt{kron}).

\textbf{Do not} use \texttt{covar\_factor} as $\RB$ unless $\LB\LB^\top = B$ has been verified;
in standard GPyTorch \texttt{IndexKernel} it does not (diagonal task variance is separate).
Build $\RB$ from \texttt{covar\_matrix} via Cholesky or symmetric eigendecomposition.
\end{diffbox}

\begin{samebox}
This is the multitask analogue of $\Sig = \noise I_n + \Ph\Ph^\top$ in the scalar note:
replace $\Ph \in \mathbb{R}^{n \times m}$ by $\Om \in \mathbb{R}^{(nT) \times (mT)}$
and noise dimension $n$ by $nT$.
\end{samebox}

\subsection{Example: TOA ($T=2$, $n=10{,}000$, $D=256$)}

\begin{center}
\renewcommand{\arraystretch}{1.2}
\begin{tabular}{@{}ll@{}}
\toprule
Quantity & Value \\
\midrule
Dense covariance size & $nT \times nT = 20{,}000 \times 20{,}000$ \\
$mT$ (joint feature width) & $512 \times 2 = 1{,}024$ \\
Woodbury middle matrix & $1{,}024 \times 1{,}024$ \\
Benefit & $1{,}024^3 \ll 20{,}000^3$ \\
\bottomrule
\end{tabular}
\end{center}

\subsection{When Woodbury helps}

Extend \texttt{\_check\_woodbury\_rank}: warn when $mT \ge nT$.
Woodbury is advantageous when $mT \ll nT$.

%=============================================================================
\section{Extension of scalar Woodbury}
%=============================================================================

This section parallels \texttt{rff\_woodbury\_derivation.tex} Sections~4--6.
Read that document for term-by-term decomposition of $(\noise I)^{-1}$ vs Cholesky of $M$.

\subsection{MVP: per-task diagonal noise (\texttt{rank\_likelihood}$=0$)}

\textbf{Phase~1} matches \texttt{MTGPR} with \texttt{rank\_likelihood}$=0$:
independent Gaussian noise per task,
\begin{equation}
  \Lambda = I_n \otimes D_\noise,
  \qquad
  D_\noise = \mathrm{diag}(\noise_1,\ldots,\noise_T)
  \quad\text{(GPyTorch \texttt{task\_noises}).}
\end{equation}
The joint covariance is
\begin{equation}
  \label{eq:sigma-mvp}
  \Sig = \Lambda + \Om\Om^\top.
\end{equation}

Woodbury inverse (structured $\Lambda$, $\Om$ replaces $\Ph$, $nT$ replaces $n$):
\begin{equation}
  \label{eq:woodbury-mt}
  \Sig^{-1}
  = \Lambda^{-1}
  - \Lambda^{-1}\,\Om\,
    \underbrace{\Bigl(I_{mT} + \Om^\top \Lambda^{-1} \Om\Bigr)^{-1}}_{M^{-1}}
    \Om^\top \Lambda^{-1}.
\end{equation}
Because $\Lambda = I_n \otimes D_\noise$, we have $\Lambda^{-1} = I_n \otimes D_\noise^{-1}$
and $\Lambda^{-1}\Om$ is cheap (scale each block of $\Om$ by the matching $\noise_t^{-1}$).

Middle matrix:
\begin{equation}
  M = I_{mT} + \Om^\top \Lambda^{-1} \Om
  \;\in\; \mathbb{R}^{(mT) \times (mT)}.
\end{equation}

\begin{samebox}
\textbf{Special case.} If all $\noise_t \equiv \noise$, then $\Lambda = \noise I_{nT}$ and
$M = I_{mT} + \Om^\top\Om / \noise$, recovering the scalar Woodbury template in
\texttt{rff\_woodbury\_derivation.tex}.
\end{samebox}

\begin{samebox}
\textbf{Reuse existing helpers} in \texttt{gpplus/utils/rff\_utils.py}:
\texttt{woodbury\_factor}, \texttt{woodbury\_solve}, \texttt{woodbury\_log\_det},
\texttt{woodbury\_marginal\_log\_likelihood} --- call them with $\Om$, structured $\Lambda$,
and vectorized $y \in \mathbb{R}^{nT}$ (row-wise \texttt{train\_y.reshape(-1)}).
Only \textbf{building} $\Om = \Ph \otimes \RB$ from spatial features and the PSD task factor
is new.
\end{samebox}

\subsection{Marginal log-likelihood}

For centered $\tilde{y} = \mathrm{vec}(y) - \mathrm{vec}(\mu)$:
\begin{equation}
  \log p(y \mid X)
  = -\tfrac{1}{2}\,\tilde{y}^\top \Sig^{-1} \tilde{y}
  - \tfrac{1}{2}\log|\Sig|
  - \tfrac{nT}{2}\log(2\pi),
\end{equation}
with
\begin{equation}
  \log|\Sig| = \log|\Lambda| + \log|M|,
  \qquad
  \log|\Lambda| = n \sum_{t=1}^{T} \log \noise_t.
\end{equation}

\subsection{Predictive mean at test points}

For test inputs $X_\ast \in \mathbb{R}^{n_\ast \times d}$, build $\Ph_\ast$ and
$\Om_\ast = \Ph_\ast \otimes \RB$.
Let $\alpha = \Sig^{-1}\tilde{y}$ (from Woodbury).
The latent cross-covariance between train and test is
$K_{x,\ast} \otimes B \approx \Om_\ast \Om^\top$ (not $\Om_\ast \Om_\ast^\top$).
Posterior mean for vectorized targets:
\begin{equation}
  \mathbb{E}[\mathrm{vec}(y_\ast)]
  = \Om_\ast \Om^\top \alpha + \mathrm{vec}(\mu_\ast),
\end{equation}
then reshape to $(n_\ast, T)$.
Per-task marginal variances use the diagonal of the predictive covariance
(\texttt{woodbury\_predictive\_var\_diag} generalized, or full $(T\times T)$ blocks per point).

\subsection{Phase~2: correlated task noise ($r_\ell > 0$)}

When \texttt{MultitaskGaussianLikelihood} has $r_\ell > 0$, $\Lambda$ has off-diagonal
task structure (low-rank task noise on top of per-task variances).
Options (when $\Lambda$ is not diagonal):
\begin{enumerate}[nosep]
  \item Whitening: $\Sig = \Lambda^{1/2}\bigl(I + \Lambda^{-1/2}\Om\Om^\top\Lambda^{-1/2}\bigr)\Lambda^{1/2}$
        and Woodbury on the whitened system;
  \item Low-rank correction if $\Lambda = \Lambda_0 + U U^\top$ with diagonal $\Lambda_0$.
\end{enumerate}
Implement Phase~1 (diagonal $\Lambda$) first; defer full $r_\ell > 0$ parity to Phase~2.

%=============================================================================
\section{GPPlus code changes}
%=============================================================================

\subsection{New and modified files}

\begin{center}
\renewcommand{\arraystretch}{1.25}
\begin{tabular}{@{}llp{0.42\textwidth}@{}}
\toprule
\textbf{File} & \textbf{Action} & \textbf{Role} \\
\midrule
\texttt{gpplus/models/rff\_mtgpr.py} & \textbf{new} & \texttt{RFFMTGPR} model \\
\texttt{gpplus/training/rff\_mt\_mll.py} & \textbf{new} & \texttt{RFFMTWoodburyMarginalLogLikelihood} \\
\texttt{gpplus/utils/rff\_utils.py} & extend & \texttt{build\_icm\_joint\_features}, MT predict/MLL wrappers \\
\texttt{gpplus/training/eval.py} & extend & \texttt{evaluate\_rff\_mt\_gp\_model} \\
\texttt{gpplus/training/parameter\_initializer.py} & extend & \texttt{RFFMTParameterInitializer} \\
\texttt{gpplus/training/training\_single\_run.py} & modify & MLL bypass for MT Woodbury \\
\texttt{gpplus/training/training\_metrics.py} & modify & validation metrics for $(n,T)$ targets \\
\texttt{gpplus/models/\_\_init\_\_.py} & modify & export \texttt{RFFMTGPR} \\
\bottomrule
\end{tabular}
\end{center}

\subsection{Math $\to$ code mapping}

\begin{center}
\renewcommand{\arraystretch}{1.2}
\begin{tabular}{@{}ll@{}}
\toprule
\textbf{Math} & \textbf{GPPlus (new)} \\
\midrule
$\Ph(x)$ & \texttt{RFFMTGPR.scaled\_spatial\_features(x)} \\
$\RB$ (PSD factor of $B$) & \texttt{task\_psd\_factor()} from \texttt{covar\_module.task\_covar.covar\_matrix} \\
$\Om = \Ph \otimes \RB$ & \texttt{RFFMTGPR.joint\_features(x)} \\
Woodbury MLL & \texttt{RFFMTWoodburyMarginalLogLikelihood} \\
Predict $(n_\ast, T)$ & \texttt{evaluate\_rff\_mt\_gp\_model} \\
ORF/RFF/SORF & \texttt{RFFKernel(rff\_sampling=...)} --- \textbf{unchanged} \\
\bottomrule
\end{tabular}
\end{center}

\subsection{\texttt{RFFMTGPR} sketch}

\begin{planbox}
\begin{verbatim}
class RFFMTGPR(gpytorch.models.ExactGP):
    # train_y: (n, T)
    # MultitaskMean, MultitaskGaussianLikelihood
    # MultitaskKernel(LogScaleKernel(RFFKernel(...)))
    # rff_sampling in {"rff", "orf", "sorf"}

    def scaled_spatial_features(self, x) -> Phi  # (n, m)
    def task_psd_factor(self) -> R_B             # (T, T); B = R_B @ R_B.T
    def joint_features(self, x) -> Omega         # (n*T, m*T); torch.kron(Phi, R_B)
    def predict(self, test_x) -> mean, lower, upper  # each (n_test, T)
    def invalidate_feature_cache(self)           # same as RFFGPR
\end{verbatim}
\end{planbox}

\subsection{Trainer routing}

In \texttt{training\_single\_run.py}, extend the Woodbury bypass:
\begin{planbox}
\begin{verbatim}
if isinstance(mll, (RFFWoodburyMarginalLogLikelihood,
                    RFFMTWoodburyMarginalLogLikelihood)):
    return -mll(None, train_y)
\end{verbatim}
\end{planbox}
Avoid dense \texttt{ExactGP} forward for large $nT$.

\subsection{What does \emph{not} change}

\begin{samebox}
\begin{itemize}[nosep]
  \item \texttt{RFFKernel}, \texttt{init\_rbf\_weights}, ORF/SORF sampling in \texttt{rff\_utils.py}
  \item Existing \texttt{RFFGPR} (scalar) and \texttt{MTGPR} (dense, small $n$)
  \item Independent $T$ scalar \texttt{RFFGPR} models (valid baseline)
\end{itemize}
\end{samebox}

%=============================================================================
\section{Implementation phases}
%=============================================================================

\subsection{Phase 1 --- MVP (joint multitask Woodbury)}

\begin{enumerate}
  \item $r_k \ge 0$ (any \texttt{rank\_kernel}), $r_\ell = 0$ (per-task diagonal noise).
  \item \texttt{RFFMTGPR} + \texttt{RFFMTWoodburyMarginalLogLikelihood} + \texttt{evaluate\_rff\_mt\_gp\_model}.
  \item \texttt{build\_icm\_joint\_features} --- \texttt{torch.kron(Phi, R\_B)} with full $B$;
        vec-order and NLL test vs dense \texttt{MTGPR} on small $n,T,m$.
  \item Refactor \texttt{toa\_*\_base.py} from 2$\times$ \texttt{RFFGPR} loop to single \texttt{RFFMTGPR}.
\end{enumerate}

\subsection{Phase 2 --- Parity with \texttt{MTGPR}}

\begin{enumerate}
  \item $r_\ell > 0$ (correlated task noise); Woodbury-compatible general $\Lambda$.
  \item \texttt{RFFMTParameterInitializer} (RFF draws + multitask hyperparameters).
  \item Multitask validation metrics in \texttt{compute\_validation\_metrics}.
  \item ORF/SORF smoke tests on \texttt{MT1\_analytic\_1d\_GP} benchmark.
\end{enumerate}

\subsection{Phase 3 --- Experiments and docs}

\begin{enumerate}
  \item Register in \texttt{run\_*\_experiments.py} batch runners.
  \item Per-task + aggregate metrics JSON (reuse \texttt{compute\_per\_task\_metrics} from TOA).
  \item Cross-link this plan with \texttt{rff\_woodbury\_derivation.tex} in repo README.
\end{enumerate}

%=============================================================================
\section{Risks and tests}
%=============================================================================

\begin{enumerate}
  \item \textbf{Vec ordering.} $\Om = \Ph \otimes \RB$ must match GPyTorch
        $\Sig_{\mathrm{latent}} = K_x \otimes B$ with $\mathrm{vec}(Y)=\texttt{train\_y.reshape(-1)}$
        (Section~\ref{sec:mtgpr}).
        Test: tiny data, compare NLL and predictions to dense \texttt{MTGPR}.
  \item \textbf{Task matrix $B$.} Use \texttt{covar\_matrix.to\_dense()}, not \texttt{covar\_factor}
        alone; verify $\RB\RB^\top = B$ before trusting $\Om\Om^\top$.
  \item \textbf{Likelihood noise.} Phase~1: diagonal $\Lambda = I_n \otimes D_\noise$;
        Phase~2: low-rank \texttt{MultitaskGaussianLikelihood} ($r_\ell > 0$) needs whitening.
  \item \textbf{ARD in high $d$.} Same as scalar RFF; $m$ depends on $D$, not $d$.
  \item \textbf{Woodbury conditioning.} Warn when $mT \ge nT$.
  \item \textbf{Cross-task predictive correlation.} MVP can report per-task RMSE;
        full $(T\times T)$ predictive cov optional for calibrated NIS.
\end{enumerate}

%=============================================================================
\section{Complexity summary}
%=============================================================================

\begin{center}
\renewcommand{\arraystretch}{1.3}
\begin{tabular}{@{}ll@{}}
\toprule
\textbf{Method} & \textbf{Training / inference cost (dominant)} \\
\midrule
Dense \texttt{MTGPR} & $O(n^3 T^3)$ Cholesky on $nT \times nT$ \\
$T$ independent \texttt{RFFGPR} & $T \cdot O(n m^2)$; no task correlation \\
\textbf{Proposed \texttt{RFFMTGPR}} & $O\bigl((mT)^3 + nT \cdot mT\bigr)$ \\
\bottomrule
\end{tabular}
\end{center}

%=============================================================================
\section{References}
%=============================================================================

\begin{thebibliography}{99}

\bibitem[Alvarez et al.(2012)]{alvarez2012kernels}
M.~\'Alvarez, L.~Rosasco, and N.~D. Lawrence.
\newblock Kernels for vector-valued functions: a review.
\newblock \emph{Foundations and Trends in Machine Learning}, 4(3):195--266, 2012.
\newblock \url{https://arxiv.org/abs/1106.6242}.

\bibitem[Bonilla et al.(2008)]{bonilla2008mtgp}
E.~V. Bonilla, K.~M. Chai, and C.~K.~I. Williams.
\newblock Multi-task Gaussian process prediction.
\newblock In \emph{Advances in Neural Information Processing Systems 20}, 2008.
\newblock \url{https://proceedings.neurips.cc/paper/2007/file/66368270ffd51418ec58bd793f2d9b1b-Paper.pdf}.

\bibitem[Cutajar et al.(2017)]{cutajar2017dgp}
K.~Cutajar, E.~V. Bonilla, P.~Michiardi, and M.~Filippone.
\newblock Random feature expansions for deep Gaussian processes.
\newblock In \emph{Proceedings of ICML}, volume~70, pages 884--893, 2017.
\newblock \url{https://proceedings.mlr.press/v70/cutajar17a.html}.

\bibitem[Gardner et al.(2018)]{gardner2018gpytorch}
J.~R. Gardner, G.~Pleiss, D.~Bindel, K.~Q.~Weinberger, and A.~G. Wilson.
\newblock GPyTorch: Blackbox matrix-matrix Gaussian process inference with GPU acceleration.
\newblock In \emph{Advances in Neural Information Processing Systems 31}, 2018.
\newblock \url{https://github.com/cornellius-gp/gpytorch}.

\bibitem[Leroy et al.(2023)]{leroy2023lmc}
V.~Leroy, F.~Gauthier, C.~Rasmussen, V.~Picheny, N.~Dutertre, and Y.~Richet.
\newblock Exact and general decoupled solutions of the linear model of coregionalization.
\newblock \emph{arXiv:2310.12032}, 2023.
\newblock \url{https://arxiv.org/abs/2310.12032}.

\bibitem[Lin et al.(2025)]{lin2025latentkronecker}
A.~Lin, S.~S. Wilson, and A.~G. Wilson.
\newblock Scalable Gaussian processes with latent Kronecker structure.
\newblock In \emph{Proceedings of ICML}, volume~267, 2025.
\newblock \url{https://proceedings.mlr.press/v267/lin25a.html}.

\bibitem[Potapczynski et al.(2021)]{potapczynski2021rfftrunc}
M.~Potapczynski, G.~Pleiss, and A.~G. Wilson.
\newblock Bias-free scalable Gaussian processes via randomized truncations.
\newblock In \emph{Proceedings of ICML}, volume~139, pages 8430--8440, 2021.
\newblock \url{https://arxiv.org/abs/2102.06695}.

\bibitem[Rahimi and Recht(2007)]{rahimi2007rff}
A.~Rahimi and B.~Recht.
\newblock Random features for large-scale kernel machines.
\newblock In \emph{Advances in Neural Information Processing Systems 20}, 2007.
\newblock \url{https://proceedings.neurips.cc/paper/2007/file/013a006f03dbc5392effeb8f18fda755-Paper.pdf}.

\bibitem[Sevilla-Salinas et al.(2023)]{sevilla2023rffblr}
F.~Sevilla-Salinas, J.~Navarro-Moreno, and C.~Mart\'inez-Ram\'on.
\newblock Bayesian learning of feature spaces for multitask regression.
\newblock \emph{arXiv:2209.03028}, 2023.
\newblock \url{https://arxiv.org/abs/2209.03028}.

\bibitem[Wilson et al.(2015)]{wilson2015kronecker}
A.~G. Wilson, C.~Dann, and H.~Nickisch.
\newblock Fast Kronecker inference in Gaussian processes with non-Gaussian likelihoods.
\newblock In \emph{Proceedings of ICML}, volume~37, pages 1864--1872, 2015.
\newblock \url{https://www.cs.cmu.edu/~neill/papers/icml15.pdf}.

\end{thebibliography}

%=============================================================================
\appendix
\section{File checklist}
%=============================================================================

\begin{itemize}
  \item[\checkmark] \texttt{docs/overleaf/rff\_woodbury\_derivation.tex} --- scalar Woodbury reference
  \item[\ ] \texttt{docs/overleaf/rff\_multitask\_gp\_plan.tex} --- this document
  \item[\ ] \texttt{gpplus/models/rff\_mtgpr.py}
  \item[\ ] \texttt{gpplus/training/rff\_mt\_mll.py}
  \item[\ ] \texttt{gpplus/utils/rff\_utils.py} --- \texttt{build\_icm\_joint\_features}, MT wrappers
  \item[\ ] \texttt{gpplus/training/eval.py} --- \texttt{evaluate\_rff\_mt\_gp\_model}
  \item[\ ] \texttt{gpplus/training/parameter\_initializer.py} --- \texttt{RFFMTParameterInitializer}
  \item[\ ] \texttt{experiments\_ORF/toa\_orf\_base.py} --- joint model (replace 2$\times$ loop)
  \item[\ ] \texttt{experiments\_RFF/toa\_rff\_base.py}
  \item[\ ] \texttt{experiments\_SORF/toa\_sorf\_base.py}
\end{itemize}

\vspace{1em}
\noindent
\textbf{Overleaf:} upload \texttt{docs/overleaf/}, set main file to \texttt{rff\_multitask\_gp\_plan.tex}, compile with pdfLaTeX.

\end{document}
